% =====================================================================
%  REPLACEMENT for the dictionary table at main.tex lines 113-124.
%  Drafted 2026-08-25. Label preserved: tab:entanglement-hopf-particles.
%
%  WHAT CHANGED AND WHY.
%  (1) Bosons and fermions now occupy separate columns, so gluons have
%      a home. They were missing because the old "Particles" column
%      held whichever species the row's SLOCC class named, and the
%      GHZ row named quarks.
%  (2) The rows are now keyed to the Cayley-Dickson SECTOR rather than
%      to the SLOCC class alone. This is forced by (1): the moment
%      bosons get their own column, a row must be a place where both
%      species can be listed, and the sector is that place. It also
%      repairs the category error flagged in
%      terminology_reconciliation_entanglement_dictionary_2026-06-08.md
%      section 5, item 4 -- the old table read as the LEVEL dictionary
%      (rung -> gauge group) in rows 1-2 and the TOPOLOGY dictionary
%      (SLOCC class -> particle) in rows 3-4.
%  (3) The gauge-group column is now explicit and uses the level
%      reading throughout, per the same note's item 5.
%
%  READING OF THE BOSON COLUMN. A gauge boson is a bipartite
%  (Bell-class) EDGE; its gauge group is the group of the sector the
%  edge connects. So every entry in the boson column is Bell-class in
%  the topology sense, including the gluons, which are bipartite colour
%  links at the octonionic rung -- NOT at the quaternionic rung. This
%  is why the boson column carries a footnote: without it, listing
%  gluons on the GHZ row re-imports the "Bell = SU(2)" misreading the
%  reconciliation note exists to prevent.
%
%  STATUS. The fermion column is PROVED (GHZ -> quark colour triplet;
%  W -> lepton singlet). The boson column is STRUCTURAL: per section 6
%  of the reconciliation note, two calculations are owed -- the four
%  electroweak bosons as a C-tensor-H bipartite quartet, and the eight
%  gluons as the SU(3) adjoint at the octonionic rung. The caption
%  carries that flag explicitly; do not drop it before those are
%  discharged.
%
%  THE HIGGS is deliberately absent from the table and remains in the
%  prose immediately following, since it is not an entanglement class
%  but the bundle's nontriviality promoted to a field. With a boson
%  column now present a reader will look for it, so the following
%  sentence should be adjusted to say so explicitly.
%
%  REQUIRES: booktabs. Set in \footnotesize to fit a4 one-column text
%  width; at 11pt \small it overruns by ~10pt. If your text block is
%  narrower still,
%  switch table -> table* (two-column layouts) or wrap in
%  \resizebox{\textwidth}{!}{...}.
% =====================================================================

\begin{table}[h]
  \centering
  \footnotesize
  \begin{tabular}{@{}lllll@{}}
    \toprule
    \textbf{Sector} & \textbf{Hopf map} & \textbf{Gauge group}
      & \textbf{Gauge bosons}\footnotemark{} & \textbf{Fermions} \\
    \midrule
    $\mathbb{C}$ (single qubit)
      & $S^{1} \to S^{3} \to S^{2}$
      & $U(1)$
      & $B$
      & --- \\
    $\mathbb{H}$ (Bell)
      & $S^{3} \to S^{7} \to S^{4}$
      & $SU(2)$
      & $W^{\pm}$, $Z$, $\gamma$
      & --- \\
    $\mathbb{O}$, $\mathbb{C}$ (W)
      & $S^{7} \to S^{15} \to S^{8}$
      & singlet
      & ---
      & Leptons \\
    $\mathbb{O}$, $\mathbb{C}^{3}$ (GHZ)
      & $S^{7} \to S^{15} \to S^{8}$
      & $SU(3)$
      & 8 gluons
      & Quarks \\
    \bottomrule
  \end{tabular}
  \ \caption{Cayley--Dickson sectors, Hopf fibrations, and the
    corresponding Standard Model species. The $\mathbb{C}$-sector boson
    is the hypercharge $B$; the photon is the post-mixing combination,
    listed with the electroweak quartet. The two octonionic rows share
    one bundle and split under the computability decomposition
    $\mathbb{O} \cong \mathbb{C} \oplus \mathbb{C}^{3}$, and that split
    is visible in the root system of the octonionic automorphism group
    $G_{2} = \operatorname{Aut}(\mathbb{O})$, whose roots come in two
    lengths. The eight long roots close into a subalgebra, and that
    subalgebra is the colour $\mathfrak{su}(3)$: the gluons are those
    eight. The six short roots do not close, forming instead the coset
    that carries the colour triplet, and it is this failure to close
    that confines the quarks, since a single colour direction cannot be
    assembled into a self-contained subalgebra. Dashes are structural,
    not omissions: matter appears only at the octonionic rung, because
    a fermion is a node carrying internal Cayley--Dickson structure and
    that structure is complete only at three qubits; and the
    $\mathbb{C}$ sector carries no colour edge, which is why leptons do
    not feel the strong force. The fermion column is established
    (Appendix~\ref{app:ghz-w-proofs}); the gauge-boson column is a
    structural identification, with the electroweak quartet and the
    gluon octet as calculations still owed.}
  \label{tab:entanglement-hopf-particles}
\end{table}
\footnotetext{Every gauge boson is a bipartite (Bell-class) edge, and
  its gauge group is that of the sector the edge connects. The gluons
  are therefore Bell-class \emph{in topology} while living at the
  octonionic rung: ``Bell'' names the two-party connector here, not the
  quaternionic level. The coincidence of the two readings in row two,
  where the Bell class and the $SU(2)$ rung genuinely agree, is special
  to the weak sector.}
